\section{Теоретическая часть}

Антенна --- устройство, предназначенное для излучения или приёма электромагнитных волн. Одна из важнейших её функций состоит в формировании излучения с определёнными направленными свойствами. Основными характеристиками направленности являются диаграмма направленности (ДН), коэффициент направленного действия (КНД) и коэффициент усиления (КУ).

Диаграмма направленности по мощности --- угловое распределение мощности излучения в единицу телесного угла:
\[
P(\theta, \varphi) = r^{2} S_{r}(\theta, \varphi),
\]
где $S_{r}$ --- радиальная компонента вектора Пойнтинга на большом расстоянии $r$ от антенны. ДН окончательно формируется в зоне Фраунгофера, определяемой соотношением
\begin{align} \label{eq:th:1}
	r \gg \frac{l^{2}}{{\lambda^{\text{т}}_{\text{с}}}}.
\end{align}

\begin{figure}[H]
	\centering
	\begin{minipage}{0.45\textwidth}
		\centering
		\hspace{0.075\linewidth}
		\includegraphics[width=0.6\linewidth]{1a}
		\subcaption{ }
		\label{fig:1a}
	\end{minipage}
	\begin{minipage}{0.45\textwidth}
		\centering
		\includegraphics[width=0.9\linewidth]{1b}
		\subcaption{ }
		\label{fig:1b}
	\end{minipage}	
	\caption{Диаграмма направленности}
	\label{fig:1}
\end{figure}

\textit{Коэффициент направленного действия} $D$ характеризует выигрыш по мощности в направлении максимального излучения:
\begin{align} \label{eq:th:2}
	D = \frac{4\pi P(\theta_m, \varphi_m)}{\int\limits_{0}^{2\pi} \int\limits_{0}^{\pi} P(\theta, \varphi) \sin \theta \, d\theta d\varphi}.
\end{align}

\textit{Коэффициент усиления} $G$ определяется как произведение КНД на КПД антенны $\eta$:
\begin{align} \label{eq:th:3}
	G = D\eta,
\end{align}
где
\begin{align} \label{eq:th:4}
	\eta = \frac{P_{\text{изл}}}{P_{\text{подв}}} = \frac{\int\limits_{0}^{2\pi} \int\limits_{0}^{\pi} P(\theta, \varphi) \sin \theta \, d\theta d\varphi}{P_{\text{подв}}}.
\end{align}

В силу принципа взаимности ДН и КНД антенны в режимах передачи и приёма совпадают.

\textit{Эффективная площадь приёма} $A$ определяется как отношение полной принимаемой мощности $P_{\text{пр}}$ к плотности потока падающего излучения $S_{\text{п}}$:
\begin{align} \label{eq:th:5}
	A = \frac{P_{\text{пр}}}{S_{\text{п}}}.
\end{align}

Величины $A$ и $D$ связаны соотношением
\begin{align} \label{eq:th:6}
	A = \frac{\bigl(\lambda^{\text{т}}_{\text{с}}\bigr)^2}{4\pi} D.
\end{align}

Цель работы --- экспериментальное определение КНД пирамидальной рупорной антенны зеркальным методом (методом Парселла) и сравнение с теоретическим расчётом. Метод опирается на идеально отражающую плоскую поверхность в зоне Фраунгофера, ориентированную параллельно излучающей апертуре (см.~\cref{fig:2}).

\begin{figure}[H]
	\centering
	\includegraphics[width=0.6\linewidth]{2}
	\caption{Метод изображений}
	\label{fig:2}
\end{figure}

Согласно методу изображений, отыскание отражённого поля сводится к нахождению поля от зеркальной антенны. В результате пересчёта мощность, принимаемая антенной, равна $P_{\text{пр}} = ADP_{\text{изл}}/(16\pi X^2)$. С учётом \cref{eq:th:6} получаем
\begin{align} \label{eq:th:7}
	\frac{P_{\text{пр}}}{P_{\text{изл}}} = \frac{D^2\bigl(\lambda^{\text{т}}_{\text{с}}\bigr)^2}{64\pi^2 X^2},
\end{align}
откуда
\begin{align} \label{eq:th:8}
	D = \frac{8\pi X}{{\lambda^{\text{т}}_{\text{с}}}} \sqrt{\frac{P_{\text{пр}}}{P_{\text{изл}}}}.
\end{align}

\begin{figure}[t]
	\centering
	\includegraphics[width=0.7\linewidth]{3}
	\caption{Блок-схема установки: 1 --- генератор, 2 --- измерительная линия, 3 --- индикатор, 4 --- согласующее устройство, 5 --- рупорная антенна, 6 --- поглощающий щит, 7 --- отражающий щит}
	\label{fig:3}
\end{figure}

Измерительная установка включает генератор СВЧ диапазона (${\lambda^{\text{т}}_{\text{с}}} \approx 3$ см), волноводный тракт с измерительной линией и индикатором, пирамидальный рупор, отражающий щит и щит с поглощающим покрытием (см.~\cref{fig:3}). Отражающий щит располагается в зоне Фраунгофера $X \gg l_{1,2}^2/{\lambda^{\text{т}}_{\text{с}}}$ и имеет размеры $L_{1,2} > X\Delta\theta_{1,2} \approx X{\lambda^{\text{т}}_{\text{с}}}/l_{1,2}$. Установка позволяет менять расстояние $X + \Delta X$ в пределах $\Delta X \sim 100$ см. В измерительной линии используется квадратичный детектор, поэтому показания индикатора пропорциональны $|E|^2$.

В несогласованном режиме поле на оси волновода, нормированное на амплитуду падающей волны, записывается в виде
\begin{align} \label{eq:th:11}
	E = e^{-ihx} + \Gamma_k e^{i\varphi_k} e^{ihx} + \Gamma e^{i\varphi} e^{ihx},
\end{align}
где $x$ --- координата от конца волноводного тракта (см.~\cref{fig:4}); $h$ --- постоянная распространения волны в волноводе; $\Gamma_k e^{i \varphi_k}$ --- коэффициент отражения от конца тракта; $\Gamma e^{i \varphi}$ --- коэффициент отражения, обусловленный отражающим щитом. Смещение антенны на $\Delta X$ приводит к появлению в последнем члене множителя $e^{ih_0 2\Delta X}$ ($h_0 = \omega/c$). В результате для $|E|^2$ имеем
\begin{multline} \label{eq:th:12}
	|E|^2 = 1 + \Gamma_{\text{к}}^2 + \Gamma^2 + 2\Gamma_{k}\Gamma\cos(\varphi - \varphi_{k} + h_0 2\Delta X) +\\+ 2\Gamma_{k}\cos(2hx + \varphi_{k}) + 2\Gamma\cos(2hx + \varphi + h_0 2\Delta X).
\end{multline}

\begin{figure}[H]
	\centering
	\includegraphics[width=0.7\linewidth]{4}
	\caption{Волноводный тракт}
	\label{fig:4}
\end{figure}

Поскольку $\Gamma_{\text{к}}$ и $\Gamma$ малы, квадратичными членами можно пренебречь:
\begin{align} \label{eq:th:13}
	|E|^2 \approx 1 + 2\Gamma_{k}\cos(2hx + \varphi_{k}) + 2\Gamma\cos(2hx + \varphi + h_0 2\Delta X).
\end{align}

КНД однозначно определяется коэффициентом $\Gamma$:
\begin{align} \label{eq:th:14}
	D = \frac{8\pi X}{{\lambda^{\text{т}}_{\text{с}}}} \Gamma.
\end{align}

\subsection{Теоретический расчет КНД}

В работе использовался рупор с размером апертуры $l_1 \times l_2 = 9.2 \times 13.4~\text{см}$, подсоединённый к волноводу размером $2.3 \times 1.0~\text{см}$.

Для расчёта КНД в зоне Фраунгофера воспользуемся интегралом Кирхгофа--Гюйгенса. Компонента поля в точке $P$ выражается через интеграл Кирхгофа:
\begin{align} \label{eq:th:1.1}
	\Psi(P) = \frac{1}{4\pi} \oint_S \left(\Psi \frac{\partial G}{\partial n} - G \frac{\partial \Psi}{\partial n}\right)\, dS,
\end{align}
где $G = e^{ikr}/r$ --- функция Грина свободного пространства. Поверхность $S$ проходит через апертуру антенны и удаляется на бесконечность в остальном пространстве. Рассматривая малые отклонения от центральной оси, интеграл приводится к виду
\begin{align} \label{eq:th:1.2}
	\Psi(P) = \frac{\Psi_0 \exp \{ikz_0\}}{i z_0 {\lambda^{\text{т}}_{\text{с}}}} \iint_\Sigma \exp \left\{\frac{ik}{2z_0}(x^2 + y^2)\right\}\, dxdy,
\end{align}
где $\Sigma$ --- площадь апертуры. Интегралы Френеля в первом приближении дают значение площади апертуры, поэтому
\begin{align} \label{eq:th:1.3}
	\Psi(P) \simeq \frac{\Psi_0}{i z_0 {\lambda^{\text{т}}_{\text{с}}}} S e^{ikz_0}.
\end{align}

Плотность потока энергии в максимуме:
\begin{align} \label{eq:th:1.4}
	S_{\max} = \frac{c}{8\pi} \frac{\Psi_0^2 S^2}{r^2 \bigl(\lambda^{\text{т}}_{\text{с}}\bigr)^2}.
\end{align}

Для эквивалентного изотропного источника:
\begin{align}\label{eq:th:1.5}
	S_r = \frac{c}{8\pi} \frac{\Psi_0^2 S}{r^2} \frac{1}{4\pi}.
\end{align}

Подставляя в выражение для КНД экспериментальную длину волны ${\lambda^{\text{п}}_{\text{с}}} \approx 3{,}3$~см (с округлением), получаем:
\begin{align} \label{eq:th:1.6}
	D = \frac{S_{\max}}{S_r} = \frac{S}{\bigl(\lambda^{\text{п}}_{\text{с}}\bigr)^2} 4 \pi
	= \frac{4\pi \cdot 9{,}2 \cdot 13{,}4}{3{,}3^2}
	= \frac{4\pi \cdot 123{,}28}{10{,}89}
	\approx 142{,}3.
\end{align}

Однако выражение \cref{eq:th:1.6} даёт завышенный КНД, что связано с неоднородностью распределения поля на апертуре рупора.
